\documentclass[11pt]{article}
% Shared preamble for the printable mini lessons (1.4 onward).
%   \documentclass[11pt]{article}
%   % Shared preamble for the printable mini lessons (1.4 onward).
%   \documentclass[11pt]{article}
%   % Shared preamble for the printable mini lessons (1.4 onward).
%   \documentclass[11pt]{article}
%   \input{../preamble}
% Build the key with:  pdflatex -jobname=<name>_KEY "\def\iskey{}\input{<name>}"
% Figures shared with the web page live in figures/Unit_X/<name>.tex and are
% pulled in with \fig{Unit_X/<name>}; the standalone package drops their own
% preamble, so their macros come from figures/preamble.tex, loaded here.
\input{../../figures/preamble}
\usepackage[margin=0.75in]{geometry}
\usepackage{array}
\usepackage{enumitem}
\usepackage{amsmath}
\usepackage[helvet]{sfmath}  % math in the same sans font as the text
\usepackage{standalone}
\pagestyle{empty}
\setlength{\parindent}{0pt}
\setlist[enumerate,1]{leftmargin=1.5em,itemsep=1em}
\setlist[enumerate,2]{label=(\alph*),leftmargin=2em,itemsep=2.2em}

\newif\ifkey
\ifdefined\iskey \keytrue \fi
\colorlet{keyred}{red!75!black}
\newcommand{\ans}[1]{\ifkey\leavevmode{\color{keyred}#1}\fi}
% vertical work space on the student copy, less on the key
\newcommand{\work}[1]{\ifkey\vspace{0.6em}\else\vspace{#1}\fi}
\newcommand{\fig}[1]{\input{../../figures/#1}}
% Title line: \lessontitle{Quartiles, IQR, and Box Plots}
\newcommand{\lessontitle}[1]{{\Large\bfseries #1\ifkey\ \textcolor{keyred}{KEY}\fi}\hfill Name \underline{\hspace{2in}}\par\vspace{0.6em}}
\newcommand{\partner}{\par\vspace{0.4em}Work with a partner. Compare answers when you finish each problem.\par}

% Build the key with:  pdflatex -jobname=<name>_KEY "\def\iskey{}\input{<name>}"
% Figures shared with the web page live in figures/Unit_X/<name>.tex and are
% pulled in with \fig{Unit_X/<name>}; the standalone package drops their own
% preamble, so their macros come from figures/preamble.tex, loaded here.
% Shared preamble for web figures.  Each figure is a standalone .tex that
% \input's this file; scripts/make-figures.sh turns it into an SVG.
\usepackage[T1]{fontenc}
\usepackage{helvet}
\renewcommand{\familydefault}{\sfdefault}
\usepackage{tikz}
\usetikzlibrary{decorations.pathreplacing}
% axis numbers in the text font, not math italic
\pgfkeys{/pgf/number format/.cd,assume math mode=true,1000 sep={}}

% dot plot: \dotplot{min}{max}{value/count,...}{axis label}
\newcommand{\dotplot}[4]{%
\begin{tikzpicture}[x=0.9cm,y=0.33cm]
  \draw[thick] (#1-0.5,0) -- (#2+0.5,0);
  \foreach \t in {#1,...,#2} {\draw (\t,0.25) -- (\t,-0.25) node[below] {\t};}
  \foreach \v/\c in {#3} {\ifnum\c>0 \foreach \k in {1,...,\c} {\fill (\v,\k*1.15) circle (0.16cm);}\fi}
  \node at ({(#1+#2)/2},-2.3) {#4};
\end{tikzpicture}}

% shape sketch histogram (12 bins): \shapehist{heights}{label}
\newcommand{\shapehist}[2]{%
  \foreach \h [count=\i from 0] in {#1} \draw[fill=gray!30] (\i,0) rectangle (\i+1,\h);
  \draw (-0.3,0) -- (12.3,0);
  \node[font=\bfseries] at (6,-2.6) {#2};}
\newcommand{\meanmark}[1]{\fill (#1,0) -- ++(-0.3,-0.75) -- ++(0.6,0) -- cycle;}
\newcommand{\medmark}[2]{\draw[very thick,dashed] (#1,0) -- (#1,#2);}

% number line with minor ticks: \nline{min}{max}{label step}{axis label}
\newcommand{\nline}[4]{%
  \draw[thick] (#1-0.5,0) -- (#2+0.5,0);
  \foreach \t in {#1,...,#2} \draw (\t,0.12) -- (\t,-0.12);
  \pgfmathtruncatemacro{\nxt}{#1+#3}
  \foreach \t in {#1,\nxt,...,#2} \draw (\t,0.2) -- (\t,-0.2) node[below] {\t};
  \node at ({(#1+#2)/2},-0.95) {#4};}
% box plot above the line: \bplot{min}{Q1}{med}{Q3}{max}
\newcommand{\bplot}[5]{%
  \draw[thick] (#1,1.0) -- (#2,1.0);
  \draw[thick] (#4,1.0) -- (#5,1.0);
  \draw[thick] (#1,0.75) -- (#1,1.25);
  \draw[thick] (#5,0.75) -- (#5,1.25);
  \draw[thick,fill=gray!20] (#2,0.55) rectangle (#4,1.45);
  \draw[very thick] (#3,0.55) -- (#3,1.45);}

% ---- added for 1.4 onward ----
% axis on y=0: \axis{min}{max}{label step}{tick step}{axis label}
\newcommand{\axis}[5]{%
  \draw[thick] (#1,0) -- (#2,0);
  \pgfmathsetmacro{\axT}{#1+#4}\foreach \t in {#1,\axT,...,#2} \draw (\t,0.1) -- (\t,-0.1);
  \pgfmathsetmacro{\axL}{#1+#3}\foreach \t in {#1,\axL,...,#2}
    \draw (\t,0.18) -- (\t,-0.18) node[below] {\pgfmathprintnumber{\t}};
  \node at ({(#1+#2)/2},-0.95) {#5};}
% box plot centred at height y: \bplotat{y}{min}{Q1}{med}{Q3}{max}
\newcommand{\bplotat}[6]{%
  \draw[thick] (#2,#1) -- (#3,#1);
  \draw[thick] (#5,#1) -- (#6,#1);
  \draw[thick] (#2,#1-0.22) -- (#2,#1+0.22);
  \draw[thick] (#6,#1-0.22) -- (#6,#1+0.22);
  \draw[thick,fill=gray!20] (#3,#1-0.42) rectangle (#5,#1+0.42);
  \draw[very thick] (#4,#1-0.42) -- (#4,#1+0.42);}
% an outlier dot: \outlier{x}{y}
\newcommand{\outlier}[2]{\fill (#1,#2) circle (2.4pt);}

\usepackage[margin=0.75in]{geometry}
\usepackage{array}
\usepackage{enumitem}
\usepackage{amsmath}
\usepackage[helvet]{sfmath}  % math in the same sans font as the text
\usepackage{standalone}
\pagestyle{empty}
\setlength{\parindent}{0pt}
\setlist[enumerate,1]{leftmargin=1.5em,itemsep=1em}
\setlist[enumerate,2]{label=(\alph*),leftmargin=2em,itemsep=2.2em}

\newif\ifkey
\ifdefined\iskey \keytrue \fi
\colorlet{keyred}{red!75!black}
\newcommand{\ans}[1]{\ifkey\leavevmode{\color{keyred}#1}\fi}
% vertical work space on the student copy, less on the key
\newcommand{\work}[1]{\ifkey\vspace{0.6em}\else\vspace{#1}\fi}
\newcommand{\fig}[1]{\input{../../figures/#1}}
% Title line: \lessontitle{Quartiles, IQR, and Box Plots}
\newcommand{\lessontitle}[1]{{\Large\bfseries #1\ifkey\ \textcolor{keyred}{KEY}\fi}\hfill Name \underline{\hspace{2in}}\par\vspace{0.6em}}
\newcommand{\partner}{\par\vspace{0.4em}Work with a partner. Compare answers when you finish each problem.\par}

% Build the key with:  pdflatex -jobname=<name>_KEY "\def\iskey{}\input{<name>}"
% Figures shared with the web page live in figures/Unit_X/<name>.tex and are
% pulled in with \fig{Unit_X/<name>}; the standalone package drops their own
% preamble, so their macros come from figures/preamble.tex, loaded here.
% Shared preamble for web figures.  Each figure is a standalone .tex that
% \input's this file; scripts/make-figures.sh turns it into an SVG.
\usepackage[T1]{fontenc}
\usepackage{helvet}
\renewcommand{\familydefault}{\sfdefault}
\usepackage{tikz}
\usetikzlibrary{decorations.pathreplacing}
% axis numbers in the text font, not math italic
\pgfkeys{/pgf/number format/.cd,assume math mode=true,1000 sep={}}

% dot plot: \dotplot{min}{max}{value/count,...}{axis label}
\newcommand{\dotplot}[4]{%
\begin{tikzpicture}[x=0.9cm,y=0.33cm]
  \draw[thick] (#1-0.5,0) -- (#2+0.5,0);
  \foreach \t in {#1,...,#2} {\draw (\t,0.25) -- (\t,-0.25) node[below] {\t};}
  \foreach \v/\c in {#3} {\ifnum\c>0 \foreach \k in {1,...,\c} {\fill (\v,\k*1.15) circle (0.16cm);}\fi}
  \node at ({(#1+#2)/2},-2.3) {#4};
\end{tikzpicture}}

% shape sketch histogram (12 bins): \shapehist{heights}{label}
\newcommand{\shapehist}[2]{%
  \foreach \h [count=\i from 0] in {#1} \draw[fill=gray!30] (\i,0) rectangle (\i+1,\h);
  \draw (-0.3,0) -- (12.3,0);
  \node[font=\bfseries] at (6,-2.6) {#2};}
\newcommand{\meanmark}[1]{\fill (#1,0) -- ++(-0.3,-0.75) -- ++(0.6,0) -- cycle;}
\newcommand{\medmark}[2]{\draw[very thick,dashed] (#1,0) -- (#1,#2);}

% number line with minor ticks: \nline{min}{max}{label step}{axis label}
\newcommand{\nline}[4]{%
  \draw[thick] (#1-0.5,0) -- (#2+0.5,0);
  \foreach \t in {#1,...,#2} \draw (\t,0.12) -- (\t,-0.12);
  \pgfmathtruncatemacro{\nxt}{#1+#3}
  \foreach \t in {#1,\nxt,...,#2} \draw (\t,0.2) -- (\t,-0.2) node[below] {\t};
  \node at ({(#1+#2)/2},-0.95) {#4};}
% box plot above the line: \bplot{min}{Q1}{med}{Q3}{max}
\newcommand{\bplot}[5]{%
  \draw[thick] (#1,1.0) -- (#2,1.0);
  \draw[thick] (#4,1.0) -- (#5,1.0);
  \draw[thick] (#1,0.75) -- (#1,1.25);
  \draw[thick] (#5,0.75) -- (#5,1.25);
  \draw[thick,fill=gray!20] (#2,0.55) rectangle (#4,1.45);
  \draw[very thick] (#3,0.55) -- (#3,1.45);}

% ---- added for 1.4 onward ----
% axis on y=0: \axis{min}{max}{label step}{tick step}{axis label}
\newcommand{\axis}[5]{%
  \draw[thick] (#1,0) -- (#2,0);
  \pgfmathsetmacro{\axT}{#1+#4}\foreach \t in {#1,\axT,...,#2} \draw (\t,0.1) -- (\t,-0.1);
  \pgfmathsetmacro{\axL}{#1+#3}\foreach \t in {#1,\axL,...,#2}
    \draw (\t,0.18) -- (\t,-0.18) node[below] {\pgfmathprintnumber{\t}};
  \node at ({(#1+#2)/2},-0.95) {#5};}
% box plot centred at height y: \bplotat{y}{min}{Q1}{med}{Q3}{max}
\newcommand{\bplotat}[6]{%
  \draw[thick] (#2,#1) -- (#3,#1);
  \draw[thick] (#5,#1) -- (#6,#1);
  \draw[thick] (#2,#1-0.22) -- (#2,#1+0.22);
  \draw[thick] (#6,#1-0.22) -- (#6,#1+0.22);
  \draw[thick,fill=gray!20] (#3,#1-0.42) rectangle (#5,#1+0.42);
  \draw[very thick] (#4,#1-0.42) -- (#4,#1+0.42);}
% an outlier dot: \outlier{x}{y}
\newcommand{\outlier}[2]{\fill (#1,#2) circle (2.4pt);}

\usepackage[margin=0.75in]{geometry}
\usepackage{array}
\usepackage{enumitem}
\usepackage{amsmath}
\usepackage[helvet]{sfmath}  % math in the same sans font as the text
\usepackage{standalone}
\pagestyle{empty}
\setlength{\parindent}{0pt}
\setlist[enumerate,1]{leftmargin=1.5em,itemsep=1em}
\setlist[enumerate,2]{label=(\alph*),leftmargin=2em,itemsep=2.2em}

\newif\ifkey
\ifdefined\iskey \keytrue \fi
\colorlet{keyred}{red!75!black}
\newcommand{\ans}[1]{\ifkey\leavevmode{\color{keyred}#1}\fi}
% vertical work space on the student copy, less on the key
\newcommand{\work}[1]{\ifkey\vspace{0.6em}\else\vspace{#1}\fi}
\newcommand{\fig}[1]{\input{../../figures/#1}}
% Title line: \lessontitle{Quartiles, IQR, and Box Plots}
\newcommand{\lessontitle}[1]{{\Large\bfseries #1\ifkey\ \textcolor{keyred}{KEY}\fi}\hfill Name \underline{\hspace{2in}}\par\vspace{0.6em}}
\newcommand{\partner}{\par\vspace{0.4em}Work with a partner. Compare answers when you finish each problem.\par}


\begin{document}

\lessontitle{Percentiles and Z-Scores}

You scored 85 on a test. Is that good? It depends on how everyone else did.
Percentiles and z-scores tell where one value stands in a data set.

\vspace{0.6em}
\textbf{Percentile}

\vspace{0.2em}
\renewcommand{\arraystretch}{1.3}
\begin{tabular}{|>{\raggedright\arraybackslash}p{6.7in}|}\hline
The \textbf{percentile} of a value is the percent of the data that is \textbf{below} it.\par\vspace{0.3em}
\hspace*{1.5in}percentile $= \dfrac{\text{number of values below}}{\text{total number of values}} \times 100$\\[1.2em] \hline
\end{tabular}

\vspace{0.4em}
\textbf{Example:} 20 students took a test. 15 students scored below 85.

\vspace{0.2em}
$15 \div 20 = 0.75$, so a score of 85 is at the \textbf{75th percentile}. 85 is higher than 75\% of the class.

\vspace{1em}
\textbf{Z-score}

\vspace{0.2em}
\begin{tabular}{|>{\raggedright\arraybackslash}p{6.7in}|}\hline
A \textbf{z-score} tells how many standard deviations a value is above or below the mean.\par\vspace{0.3em}
\hspace*{1.2in}$z = \dfrac{x - \mu}{\sigma}$ \qquad $x$ = the value, \; $\mu$ = the mean, \; $\sigma$ = the SD\par\vspace{0.4em}
$z > 0$: above the mean. \qquad $z < 0$: below the mean. \qquad $z = 0$: equal to the mean.\\ \hline
\end{tabular}

\vspace{0.4em}
\textbf{Example:} The test has a mean of 70 and an SD of 10. What is the z-score of 85?

\vspace{0.2em}
$z = \dfrac{85 - 70}{10} = \dfrac{15}{10} = 1.5$ \qquad 85 is \textbf{1.5 standard deviations above} the mean.

\begin{center}\fig{Unit_2/2.3_z_line}\end{center}

{\small $\mu$ is the Greek letter \textbf{mu}. $\sigma$ is a small \textbf{sigma}. Do not mix it up with the capital $\Sigma$, which means add up.}

\vspace{0.8em}
\textbf{Comparing with z-scores.} Z-scores let you compare values from different data sets.

\vspace{0.3em}
\begin{tabular}{|l|c|c|c|c|}\hline
 & Score & Mean & SD & $z$ \\ \hline
Ana, math test & 85 & 70 & 10 & $\frac{85 - 70}{10} = 1.5$ \\[0.3em] \hline
Ben, history test & 80 & 74 & 3 & $\frac{80 - 74}{3} = 2$ \\[0.3em] \hline
\end{tabular}

\vspace{0.4em}
Ana's score is higher, but Ben did better \textbf{compared to his class}. His z-score is larger.

\vspace{0.8em}
\textbf{Going backward:} \quad $x = \mu + z\sigma$ \qquad
$z = -2$ on Ana's test: \; $70 + (-2)(10) = 50$

\partner

% ---------------- PROBLEMS ----------------
\newpage
\begin{enumerate}

\item Here are the quiz scores of 20 students, in order.
\begin{center}52, 58, 61, 64, 66, 68, 70, 71, 73, 75, 76, 78, 80, 81, 83, 85, 88, 90, 94, 98\end{center}
\begin{enumerate}[itemsep=1.6em]
  \item What percentile is a score of 83? \ans{\quad 14 scores are below 83. $14 \div 20 = 0.70$, the 70th percentile.}
  \item What percentile is a score of 61? \ans{\quad 2 scores are below. $2 \div 20 = 0.10$, the 10th percentile.}
  \item Which score is at the 50th percentile? \ans{\quad 50\% of 20 is 10 scores below it: 76.}
\end{enumerate}
\work{1.5em}

\item This dot plot shows how many hours 25 students spent on homework last week.
\begin{center}\fig{Unit_2/2.3_dot_plot_hours}\end{center}
\begin{enumerate}[itemsep=1.6em]
  \item What percentile is 7 hours? \ans{\quad $1 + 2 + 3 + 5 + 5 = 16$ below. $16 \div 25 = 0.64$, the 64th percentile.}
  \item What percentile is 4 hours? \ans{\quad 3 below. $3 \div 25 = 0.12$, the 12th percentile.}
\end{enumerate}
\work{1.5em}

\item A class took a quiz. The mean is 20 points and the SD is 4 points. Find the z-score of each score. Then say what it means.
\begin{enumerate}[itemsep=3.5em]
  \item 26 \ans{\quad $z = \frac{26 - 20}{4} = 1.5$. \; 1.5 SDs above the mean.}
  \item 14 \ans{\quad $z = \frac{14 - 20}{4} = -1.5$. \; 1.5 SDs below the mean.}
  \item 20 \ans{\quad $z = 0$. \; Equal to the mean.}
  \item 30 \ans{\quad $z = \frac{30 - 20}{4} = 2.5$. \; 2.5 SDs above the mean.}
\end{enumerate}

\newpage
\item Lin took two tests.

\begin{center}
\renewcommand{\arraystretch}{1.3}
\begin{tabular}{|l|c|c|c|}\hline
Test & Lin's score & Class mean & Class SD \\ \hline
Spanish & 88 & 80 & 4 \\ \hline
Science & 90 & 84 & 6 \\ \hline
\end{tabular}
\end{center}
\begin{enumerate}[itemsep=1.6em]
  \item Find Lin's z-score on each test. \ans{\quad Spanish: $\frac{88 - 80}{4} = 2$. \; Science: $\frac{90 - 84}{6} = 1$.}
  \work{2em}
  \item Lin's science score is higher. On which test did Lin do better compared to the class? Explain.

  \ans{Spanish. Lin was 2 SDs above the mean in Spanish, but only 1 SD above in science.}
\end{enumerate}
\work{2em}

\item A test has a mean of 70 and an SD of 8.
\begin{enumerate}[itemsep=1.6em]
  \item What score has a z-score of 2? \ans{\quad $70 + 2 \times 8 = 86$}
  \item What score has a z-score of $-0.5$? \ans{\quad $70 + (-0.5)(8) = 70 - 4 = 66$}
\end{enumerate}
\work{1.5em}

\item In a 100-meter race, a \textbf{lower} time is better. The team's mean time is 14.0 seconds, with an SD of 0.5 seconds.
\begin{enumerate}[itemsep=1.6em]
  \item Maria ran 13.0 seconds. Find her z-score. \ans{\quad $z = \frac{13.0 - 14.0}{0.5} = -2$}
  \item Is a negative z-score good or bad for Maria? Explain.

  \ans{Good. She is 2 SDs below the mean time, so she was much faster than most of the team.}
\end{enumerate}
\work{1.5em}

\item A student says, ``A z-score of $-1$ means the value is 1 less than the mean.'' Is the student right? Explain.

\ans{No. It means the value is 1 \textbf{standard deviation} below the mean. With an SD of 10, that is 10 less than the mean.}

\end{enumerate}

\end{document}
